\section{Introduction}\label{sec:intro}

%% ⛔ PROSA SUBSTITUIDA em 22/09/2026 pela reescrita do parecer v3, que corta a introducao
%% de 779 para ~440 palavras. ⚠️ A versao dela nao traz NENHUM numero medido; as macros
%% que o criterio 42 cobra voltam abaixo, no paragrafo do desenho, por macro e nao
%% digitadas. O mapa numero->citekey da bibliografia dela esta em mapa_citacoes.json.

Multi-video summarization selects content that represents shared and complementary information across source videos~\citep{apostolidis2021survey,li2010videommr,panda2017collaborative,truong2007video_abstraction,wu2020mvsgcn}. In a segment-based pipeline, redundancy detection groups candidate segments, and a subsequent selection criterion chooses a representative from each group. Human-strategy-based summarization makes these content-selection criteria explicit~\citep{barbieri2021tese,barbieri2021content}, providing a modular setting in which one selection component can be examined while the others remain fixed.

A criterion states the property sought; its implementation supplies the decision rule and the measurements used by that rule. Here, the criterion examined is the quality criterion, and the property it is taken to seek is technical quality: fidelity of the visual signal under impairments such as noise, blur and compression artefacts, as distinct from semantic relevance, coverage, diversity and the usefulness of a whole summary. Operations developed for those purposes are reused to rank candidates by quality without being validated for that role. Tested in that role, the evaluated operations fail systematically, and which of them is used changes what the pipeline selects.

The visual-selection literature uses colourfulness, contrast, sharpness, entropy, saliency and local-feature statistics in several roles, and names them as quality: an average of five attributes is called a \emph{quality score}~\citep{srinivas2016rank}, a weighted sum of sharpness and brightness evaluates frames on what its authors call \emph{perceptual quality}~\citep{korolkov2025keyframe}, four features feed a \emph{photo-quality} classifier~\citep{ke2006highlevel}, and one source names technical fidelity outright and implements it with BRISQUE~\citep{xu2024property}. Others enter as terms of multi-objective scores or as discard filters~\citep{cakmak2025tripss,gygli2014summe,huang2024aesthetic,jeong2016static,qaroush2025diversity,yang2024keyframe}. Each source's evidence concerns the property or system it evaluates: an attribute may discriminate photographic styles, or a filter may remove blurred frames~\citep{hasler2003colorfulness,ke2006highlevel}. Among the sources verified here, none validates the isolated reuse as a technical-fidelity score. This is a question of \emph{construct validity}~\citep{cronbach1955construct}: whether an operational measure represents the property attributed to it.

The question addressed here is whether the evaluated implementations remain suitable when used as standalone technical-quality scores. Summary-level benchmarks cannot answer it: F-scores on SumMe~\citep{gygli2014summe} and TVSum~\citep{song2015tvsum} evaluate a whole summary, and randomly generated summaries reach scores comparable to the state of the art~\citep{otani2019rethinking}, so a benchmark score does not isolate whether one criterion measures its property. The audit therefore works at the level of the component.  Dedicated no-reference image quality assessment (NR-IQA) models provide relevant comparators because recordings from different broadcasters do not share an undistorted reference~\citep{agnolucci2024qualiclip,ke2021musiq,mittal2012brisque,mittal2013niqe,zhang2023liqe}. Their suitability is tested here on the intended data and decision.

We investigate the degradation response, external opinion-score association and selection
consequences of these reuses, over \fontesDaLiteratura{} operations reimplemented from
\fontesTrabalhos{} source works plus the ITU-T spatial information index, on 800
original--transformed comparisons drawn from 20 broadcast segments, and on three public image-quality databases with synthetic (KADID-10k) and authentic (KonIQ-10k, LIVE In the Wild) distortions. The objective is to establish which of these operations respond to a technical-quality requirement, where they fail, and how much the choice among them changes what a pipeline selects. The study contributes (i) a source-linked inventory of operations and decision roles, each fixed by a verbatim excerpt; (ii) a common protocol connecting controlled perturbations and external opinion databases to component substitution; (iii) an exact decomposition and empirical diagnosis of the TF-IDF-weighted SIFT+BoVW score ~\citep{barbieri2021tese}; and (iv) a controlled measurement of representative-selection sensitivity within a fixed modular testbed. HSMVideoSumm~\citep{moraes2024hsmvideosumm} serves as the modular testbed, and its own visual-word operation ~\citep{barbieri2021tese} is audited on the same footing as the others. 
